\chapter{Conclusions and Future Work}

\section{Summary of the Thesis}

This thesis investigated a 38~GHz metasurface transmitarray lens for lens-assisted multiple-input multiple-output (MIMO) communication. The central question was whether passive field focusing can provide benefits beyond received-power enhancement by changing the spatial structure of the channel. The evaluated indicators therefore included channel magnitude and link capacity as well as normalized channel coherence, singular-value balance, condition number, effective rank, and ideal MIMO capacity.

Chapter~2 developed the patch feed antenna and the $20\times20$-element planar transmitarray lens in Ansys HFSS. It formulated a common CSI representation for HFSS S-parameters and Sionna-RT frequency responses, together with raw and Frobenius-normalized MIMO metrics. The lens effect was then cross-verified numerically for a $3\times3$ focal-plane configuration using HFSS SBR+, HFSS Hybrid, and Sionna-RT. Chapter~3 extended the evaluation to two system-level studies in a larger indoor scene. The first performed a paired comparison of complete lens-assisted and without-lens $3\times3$ RX designs on matched linear and half-circular TX trajectories. The second compared complete lens-assisted $3\times3$, $5\times5$, and $9\times9$ configurations on the same linear trajectory while treating the results as configuration-level rather than order-only effects.

The results show that conclusions about a lens depend on the evaluation scale, receiver operation, and metric normalization. At the component-level $3\times3$ snapshot, the lens consistently concentrates energy into direction-matched links and reduces normalized channel coherence. At the mobile system level, however, its directional branches are highly selective, and equal weighting of aligned and non-aligned RX cases produces lower median link and constructed-MIMO performance than the more uniform without-lens baseline. Increasing MIMO order raises link capacity and total normalized constructed-MIMO capacity, but the additional dimensions are progressively less balanced and less efficiently utilized.

\section{Principal Findings}

\subsection{Antenna and Transmitarray-Lens Design}

The designed transmitarray contains $20\times20$ unit cells. The multilayer unit-cell library provides the phase states used to approximate the required focusing profile while retaining acceptable transmission around the 38~GHz design frequency. Full-wave simulation confirms that the lens produces a localized focal region and that the focus moves along a focal arc as the incidence angle changes. The imported lens-assisted RX patterns likewise show a progressive shift of the main lobe across the evaluated steering states, with moderately degraded focal sharpness and increased sidelobes at the largest angles.

The realized axial focus occurs at approximately 14~mm rather than the 30~mm design target. This discrepancy is consistent with the combined effects of incomplete and quantized phase coverage, nonuniform transmission magnitude, finite-aperture behavior, and full-wave interactions among dissimilar cells. The result demonstrates that the phase profile creates the intended focusing mechanism, but it also shows that the target focal length cannot be assumed to equal the realized focal length after physical implementation of the phase mask. The simulated focal position must therefore be used when placing focal-plane RX elements.

\subsection{Cross-Solver Channel-Model Comparison}

The three solvers agree on the principal lens-induced channel reorganization. Averaged over 37, 38, and 39~GHz, the direction-matched diagonal links increase by 8.33~dB in HFSS Hybrid, 5.33~dB in HFSS SBR+, and 7.12~dB in Sionna-RT. The corresponding mean off-diagonal changes are $-7.58$, $-7.86$, and $-7.48$~dB. All three solvers agree on the direction of the lens-induced change at 25 of the 27 evaluated link--frequency points, or 92.6\%. This agreement provides consistent qualitative evidence that the lens reinforces matched angular channels and suppresses most mismatched channels, although the link--frequency points are not statistically independent and the solvers do not predict identical absolute magnitudes.

At 38~GHz, every solver also predicts lower mean off-diagonal RX- and TX-side normalized channel coherence with the lens. HFSS SBR+ changes from 0.77/0.77 to 0.21/0.19, HFSS Hybrid from 0.39/0.56 to 0.09/0.13, and Sionna-RT from 0.30/0.34 to 0.18/0.18. Raw capacity at the 20~dB reference SNR increases from 0.60 to 1.44~bit/s/Hz in SBR+, from 0.38 to 1.11~bit/s/Hz in Hybrid, and from 0.32 to 1.03~bit/s/Hz in Sionna-RT. The consistency of these directions supports the channel-modeling approach at a qualitative level.

The spatial-mode conclusion is more solver-dependent. SBR+ improves the linear condition number from 6.25 to 2.22, effective rank from 1.61 to 2.36, and Frobenius-normalized capacity from 8.02 to 9.92~bit/s/Hz. Hybrid produces the same trend: its respective metrics change from 4.42 to 1.58, from 2.06 to 2.73, and from 9.04 to 10.46~bit/s/Hz. Sionna-RT starts from a better-conditioned without-lens channel: its linear condition number changes only from 2.06 to 1.92, effective rank decreases slightly from 2.60 to 2.56, and normalized capacity changes from 10.23 to 10.20~bit/s/Hz. The cross-solver comparison therefore supports the focusing, raw-capacity, and decorrelation trends, but it does not establish solver-independent numerical gains in normalized multiplexing capability.

\subsection{System-Level Receiver-Configuration Comparison}

The controlled $3\times3$ comparison reverses the simple expectation of a uniform lens gain. Under equal statistical weighting of all three RX configurations, the without-lens baseline has the stronger aggregate link on both trajectories. For the linear path, the lens changes median channel magnitude from $-79.74$ to $-86.75$~dB and link capacity from 5.41 to 3.57~bit/s/Hz. For the half-circular path, the corresponding changes are $-80.89$ to $-87.29$~dB and 5.08 to 3.53~bit/s/Hz. Thus, the median link-capacity reductions are 1.84~bit/s/Hz on the linear path and 1.55~bit/s/Hz on the half-circular path.

The constructed-MIMO metrics show the same capacity and effective-rank ordering, although the condition number is path dependent. Before capacity is calculated, each waypoint-level channel tensor is scaled to unit mean coefficient power over its RX, TX, and frequency dimensions; this convention removes its common amplitude scale but is distinct from the unit-Frobenius normalization used in Chapter~2. On the linear trajectory, effective rank changes from 1.202 without the lens to 1.154 with the lens, while capacity at a normalized SNR of 10~dB changes from 6.12 to 5.85~bit/s/Hz. The condition number worsens slightly from 25.01 to 25.55~dB. On the half-circular trajectory, effective rank changes from 1.348 to 1.305 and capacity from 6.71 to 6.58~bit/s/Hz, whereas the condition number improves from 25.58 to 21.56~dB. The sizeable conditioning improvement on the half-circular path therefore does not translate into a higher median effective rank or capacity, confirming that condition number alone is insufficient to establish a multiplexing advantage.

The lens nevertheless produces much stronger local angular differentiation. Median block-level decorrelation increases from 0.0909 to 0.5308 on the linear path and from 0.0638 to 0.5518 on the half-circular path. Global pooled decorrelation also favors the lens, but only slightly: it changes from 0.9784 to 0.9807 on the linear path and from 0.9744 to 0.9872 on the half-circular path. The two estimators therefore agree in direction but differ strongly in effect size because global pooling is dominated by the complete concatenated response, whereas block-level statistics characterize a typical local channel state. Moreover, these decorrelation metrics use only the center TX element and should not be equated with the singular-value behavior of the full constructed $3\times3$ matrix.

The static and mobile results together explain the system-level behavior. In the static snapshot, all three lens branches are stronger and more differentiated than their ordinal without-lens counterparts; the constructed snapshot also shows higher effective rank and capacity with the lens, despite a worse condition number. Along a trajectory, however, the angle-specific responses sweep and cross as alignment changes. The median within-waypoint capacity spread reaches 5.32~bit/s/Hz on the linear path and 4.42~bit/s/Hz on the half-circular path with the lens, compared with only 0.01~bit/s/Hz on either path without it. Equal-weight aggregation consequently penalizes the non-aligned lens branches. The present evidence therefore characterizes the lens as an angularly selective front end rather than as a source of uniform gain across all three focal-plane branches. Whether beam selection or combining can convert this selectivity into a communication benefit remains to be evaluated.

The half-circular path gives higher effective rank and constructed-MIMO capacity for both RX designs and avoids the severe conditioning peak at the midpoint of the linear path. This is descriptive because the path geometry and sampling also differ.

\subsection{Higher-Order MIMO Scaling}

For the lens-assisted order comparison, median link capacity increases from 3.57~bit/s/Hz for $3\times3$ to 4.13~bit/s/Hz for $5\times5$ and 5.19~bit/s/Hz for $9\times9$. The increase from $3\times3$ to $9\times9$ is 1.62~bit/s/Hz, or 45.4\%. This absolute link improvement is not power-normalized: the simulations assign 10~dBm per TX, so nominal total TX power rises from 14.77 to 19.54~dBm. TX aperture and RX angular sampling also change with order.

At the fixed normalized SNR of 10~dB, total constructed-MIMO capacity increases from 5.851 to 8.254 and 14.562~bit/s/Hz. The unit-mean-coefficient normalization allows the total normalized channel energy to grow with matrix dimension, so this increase should be interpreted together with the per-dimension metrics. Effective rank rises from 1.154 to 1.332 and 2.043. However, the median condition number worsens monotonically from 25.55 to 34.24 and 42.88~dB. Effective-rank fraction falls from 0.385 to 0.266 and 0.227, while capacity per dimension falls from 1.950 to 1.651 and 1.618~bit/s/Hz. The $3\times3$ to $9\times9$ transition therefore increases total normalized fixed-SNR MIMO capacity by 148.9\%, but capacity per dimension decreases by 17.0\% and effective-rank fraction decreases by 41.0\%.

All orders reach a near-rank-one bottleneck at the 9.5~m midpoint. There, the condition number is 55.09, 62.06, and 71.89~dB, and the effective rank is 1.035, 1.061, and 1.085 for $N=3$, $5$, and $9$, respectively. The shared impairment shows that higher order does not remove the geometry-related conditioning bottleneck; it raises total normalized constructed-MIMO capacity with diminishing spatial efficiency.

\section{Overall Conclusions and Contributions}

The component-level hypothesis that a lens can separate angular channels is supported. The full-wave focal analysis, diagonal channel reinforcement, off-diagonal suppression, lower normalized channel coherence, and cross-solver agreement all show that the transmitarray performs a meaningful spatial transformation. In HFSS SBR+ and Hybrid this transformation also improves singular-value balance and normalized capacity. Sionna-RT confirms the focusing and raw-capacity gains but predicts little change in normalized mode balance.

The broader hypothesis that the same lens necessarily improves a mobile multi-branch MIMO system is not supported under the operating rule evaluated in Chapter~3. With all RX configurations weighted equally, the lens-assisted $3\times3$ design has lower median link capacity, effective rank, and constructed-MIMO capacity than the without-lens baseline on both controlled paths. The lens instead provides local angle selectivity and typical-block decorrelation. Beam-selection and combining strategies must be evaluated explicitly to determine whether these properties can yield a system-level advantage while avoiding persistently non-aligned branches.

The principal contributions of this thesis are therefore:
\begin{itemize}
    \item the design and full-wave characterization of a planar $20\times20$ transmitarray lens and focal-plane RX arrangement at 38~GHz;
    \item a common raw and normalized MIMO evaluation framework for HFSS S-parameters and Sionna-RT channel responses;
    \item three-solver evidence that matched-link gain, mismatched-link suppression, raw-capacity increase, and normalized-coherence reduction are qualitatively robust, with 92.6\% agreement on the sign of the link-level lens effect;
    \item a paired mobile-trajectory result showing that component-level focusing and local decorrelation do not automatically produce higher aggregate link or multiplexing performance under equal weighting of directional RX branches;
    \item an explicit demonstration that spatial-decorrelation conclusions depend on whether normalized channel coherence is globally pooled, evaluated locally by waypoint, or summarized over typical time--frequency blocks; and
    \item a common-trajectory higher-order comparison showing growth in total normalized constructed-MIMO capacity together with worsening conditioning and declining per-dimension efficiency.
\end{itemize}

Taken together, these contributions support a conditional rather than universal conclusion. Within the evaluated simulations, the lens functions as an effective passive angular filter and focusing device. Its value for MIMO communication depends on receiver placement, alignment, channel geometry, aggregation rule, and signal-processing strategy. Array size and lens gain should therefore not be used alone as proxies for spatial-multiplexing quality.

\section{Limitations}

The conclusions are subject to the following limitations:
\begin{itemize}
    \item \textbf{Simulation-only validation}: no fabricated antenna, transmitarray, or over-the-air channel measurement is available. Cross-solver agreement supports qualitative robustness but is not experimental proof.
    \item \textbf{Solver non-equivalence}: HFSS SBR+, HFSS Hybrid, and Sionna-RT use different propagation approximations, antenna representations, and channel-extraction procedures. Their absolute magnitudes and baseline conditioning are not numerically interchangeable.
    \item \textbf{Constructed receiver matrices}: the Sionna-RT $N\times N$ matrices stack $N$ sequential $N\times1$ RX simulations. They do not model simultaneous multiport reception, mutual coupling, common hardware phase, switching overhead, or joint receiver noise.
    \item \textbf{Confounded RX comparison}: the lens-assisted and without-lens $3\times3$ cases differ in both RX placement and imported far-field pattern. The controlled trajectories compare complete receiver designs and cannot attribute the observed difference to the lens material alone.
    \item \textbf{Confounded order comparison}: TX count, nominal total power, aperture, number of RX cases, and RX angular grid change together for $N=3$, $5$, and $9$. Only the $0^\circ$ RX branch is common to all three orders.
    \item \textbf{Limited propagation sampling}: the system-level results cover one indoor scene and the available controlled trajectories without independent random seeds, confidence intervals, or a convergence sweep. The simulations use 100{,}000 path samples per source rather than the 1{,}000{,}000 recommended by the source workflow for final high-accuracy evaluation.
    \item \textbf{Metric and operating assumptions}: link capacity is an ideal Shannon bound, and constructed-MIMO capacity uses equal power at a fixed normalized SNR after scaling each waypoint-level channel to unit mean coefficient power. Neither represents coded throughput with interference, channel-estimation error, RF loss, beam-training overhead, or practical modulation and coding.
    \item \textbf{Temporal scope of trajectory metrics}: the link-budget SNR and capacity and all constructed-MIMO metrics use the first native mobility time sample at each waypoint, whereas the channel-magnitude summaries and spatial-decorrelation analysis use all 16 time samples. These metrics therefore have different temporal aggregation scopes.
    \item \textbf{Limited reliability and spatial estimators}: the 1~bit/s/Hz outage threshold is saturated at 0\% for all Chapter~3 configurations. Spatial decorrelation also depends on aggregation scope and on the order-dependent RX-pair population.
    \item \textbf{Reproducibility of the trajectory archive}: the source notebooks and derived Chapter~3 tables and figures are retained, and the packaged consistency checks pass. However, the complete complex channel tensors are not archived separately; alternative channel-level analyses therefore require rerunning the computationally expensive ray-tracing notebooks.
\end{itemize}

\section{Future Work}

The next stage should focus on converting the observed angular selectivity into a measured and operational communication benefit:
\begin{itemize}
    \item \textbf{Prototype and channel-sounder validation}: fabricate the patch antenna and transmitarray, measure the realized focal distance and angle-dependent pattern, and acquire simultaneous coherent multiport channels for the with-lens and baseline arrays.
    \item \textbf{Practical beam management}: evaluate beam selection, switching, combining, and hybrid beamforming with training overhead and alignment error. Performance should be compared with the equal-weight results to determine whether strong local branches compensate for non-aligned beams.
    \item \textbf{Factor-isolated placement experiments}: compare co-located RX elements using controlled common patterns, then vary focal displacement and imported pattern separately. This would distinguish the effects of lens transformation, physical placement, and antenna radiation pattern.
    \item \textbf{Fair order scaling}: repeat the $3\times3$, $5\times5$, and $9\times9$ study with fixed total TX power, fixed physical aperture, and a common RX angle set. Absolute capacity and per-dimension efficiency should be reported together.
    \item \textbf{Statistical and reproducible evaluation}: archive the complete complex channel-frequency-response tensor together with the source notebooks, increase ray samples, use multiple seeds and trajectories, perform convergence analysis, and report confidence intervals and informative capacity percentiles or outage thresholds.
    \item \textbf{Lens re-optimization}: expand the unit-cell phase library and include full-wave focal-position error in the design loop to reduce the difference between the 30~mm target and the approximately 14~mm realized focus. Wideband behavior and beam squint should also be evaluated.
    \item \textbf{Richer deployment scenarios}: include blockage, clutter, multiple rooms, outdoor links, orientation errors, and moving terminals. Doppler spread, coherence time, and tracking stability should be assessed together with capacity.
    \item \textbf{Multi-user and higher-order operation}: investigate interference-aware precoding, user scheduling, focal-plane beam selection, and multi-lens architectures. Orders above $N=9$ should be pursued together with strategies that preserve effective-rank fraction rather than increasing element count alone.
\end{itemize}

In conclusion, the simulated transmitarray lens demonstrates consistent passive focusing and angular channel separation at 38~GHz, but a communication-system gain is not automatic. A central design question is how to match the lens's directional response with receiver placement and adaptive signal processing. Future lens-assisted MIMO systems should therefore be optimized jointly across electromagnetic design, channel geometry, and beam-management algorithms.
